\chapter{Picard--Vessiot Rings}

For an algebraic polynomial one adjoins roots. For a linear differential
or difference system one adjoins an invertible matrix of solutions,
and the role of the splitting field is played initially by a ring.
This chapter explains why the difference ring can have zero divisors,
why no new constants appear under the stated hypotheses, and how to
recognize the resulting extension. The two theories will use the same
formalism in the next chapter.

\section{The operator and the system}

\begin{definition}[Operator rings and constants]
\label{res:PV01}
All rings in these chapters are commutative and have an identity. In
the differential case an operator is a derivation \(\partial\), an
additive map satisfying \(\partial(ab)=a\partial b+b\partial a\).
In the difference case it is a ring automorphism \(\sigma\).
An operator subring is closed under \(\partial\), or under both
\(\sigma\) and \(\sigma^{-1}\), respectively. An operator ideal is
defined by the same stability condition. An operator ring is
\emph{simple} if it is nonzero and its only operator ideals are \(0\)
and the whole ring. Its constants are the kernel of \(\partial\), or
the fixed ring of \(\sigma\), respectively, and will be denoted
\(S^{\mathrm{const}}\) for either kind of operator on \(S\).
An operator field is a field with the indicated operator.
\end{definition}

Except where explicitly stated, \(k\) is an operator field of
characteristic zero and \(C=k^{\mathrm{const}}\) is algebraically
closed. On a tensor product of extensions over \(k\) the differential
operator is the sum of the operators in the two factors and the
difference operator acts in both factors. The product rule and the
common operator on \(k\) make these formulas well defined.
Finite-type algebras over a noetherian base are noetherian
by the Hilbert basis theorem in~\cite[Section 23.11, PDF p.~191]{mg-d07}.

\begin{lemma}[Fundamental matrices and scalar equations]
\label{res:PV02}
In the differential case let \(A\in M_n(k)\) and consider
\(\partial Y=AY\). In the difference case let \(A\in GL_n(k)\)
and consider \(\sigma(Y)=AY\). An invertible solution matrix is
called fundamental. If \(Y\) is fundamental in an extension \(S\),
every solution column in \(S^n\) is uniquely of the form \(Yc\) with
\(c\in(S^{\mathrm{const}})^n\). Two fundamental matrices in \(S\)
differ on the right by an invertible constant matrix.
A monic scalar equation of order \(n\) gives such a system by taking
the column consisting of the unknown and its first \(n-1\) derivatives
or iterates. For the difference equation its order-zero coefficient
must be nonzero to give an invertible system matrix.
\end{lemma}

\begin{proof}
For a solution column \(v\), differentiating \(Y^{-1}v\) gives
\(-Y^{-1}Av+Y^{-1}Av=0\). In the difference case
\(\sigma(Y^{-1}v)=(AY)^{-1}Av=Y^{-1}v\).
The converse follows by applying the operator to \(Yc\).
Apply this to the columns of a second fundamental matrix and then
interchange the two matrices to obtain constant inverses.
For the scalar equation, the first \(n-1\) rows of its system shift the
entries of the column by one. The last row is the negative of the
coefficients in increasing order. In the difference case its determinant
is the order-zero coefficient times a sign, so the stated condition
is precisely invertibility.
\end{proof}

\section{Simplicity, constants and the total ring}

\begin{lemma}[The components of a simple operator ring]
\label{res:PV03}
A simple noetherian differential ring containing \(\mathbb Q\) is a
domain. A simple noetherian difference ring is a finite product of
domains whose factors are permuted transitively by the automorphism.
In either case the operator extends to the total ring of fractions,
obtained by inverting the non-zero-divisors. This ring is a field in
the differential case and a finite product of fields with a transitive
permutation in the difference case. A simple ring and its total ring
of fractions have the same constants, and these constants form a field.
\end{lemma}

\begin{proof}
If \(x^m=0\), applying a derivation \(m\) times shows that
\(m!(\partial x)^m\in(x)\), because the terms other than the one
with one derivative in each factor have an undifferentiated factor.
Thus the nilradical is a differential ideal in characteristic zero
when the ring contains \(\mathbb Q\). For every minimal prime \(P\),
localize at \(P\) and apply this argument to the nilradical \(PR_P\).
Contraction to the original ring shows that \(P\) is a differential
ideal. Simplicity gives \(P=0\), so the ring is a domain.

In the difference case the nilradical is a proper stable ideal, so
the ring is reduced. The automorphism permutes its finite set of
minimal primes. To recall finiteness, a minimal counterexample
among closed sets of a noetherian spectrum would split into two
smaller closed sets, each a finite union of irreducible sets.
Applied to the whole spectrum, this gives finitely many
irreducible components and hence finitely many minimal primes.
The intersection of any proper subset of the
minimal primes is nonzero: choose an omitted minimal prime \(P\)
and multiply elements, one in each prime of the subset and outside
\(P\). This also applies to the empty subset, using the product \(1\).
Simplicity therefore implies transitivity. For the minimal primes
\(P_i\), the sum of the ideals \(J_i=\bigcap_{j\ne i}P_j\) is nonzero
and is stable, so it equals the ring. Write \(1=\sum_i e_i\) with
\(e_i\in J_i\). The image of \(e_i\) modulo \(P_j\) is \(1\) for
\(i=j\) and \(0\) otherwise. The intersection of the \(P_j\) is zero,
so the \(e_i\) are orthogonal idempotents. The components \(e_iS\)
are the domains \(S/P_i\). A non-zero-divisor has a nonzero image
in each component, which proves the assertion about the total
ring of fractions.

The quotient rule extends \(\partial\), and
\(\sigma(r/s)=\sigma(r)/\sigma(s)\) extends the automorphism.
A nonzero constant of a simple ring generates a nonzero stable
ideal, hence is a unit with constant inverse. If \(q\) in its total
ring of fractions is constant, the ideal
\(\{a\in S:aq\in S\}\) is stable and contains a denominator of \(q\).
It therefore contains \(1\), so \(q\) was already a constant of \(S\).
\end{proof}

\begin{lemma}[Specialization in a finitely generated domain]
\label{res:PV04}
Suppose that \(B\) is a finitely generated domain over an infinite
field \(F\) and \(t\in B\) is transcendental over \(F\).
Then \(t-a\) is a nonunit of \(B\) for all but finitely many \(a\in F\).
Consequently a field finitely generated as an algebra over \(F\)
is finite algebraic over \(F\). If \(F\) is also algebraically
closed, a nonzero reduced finitely generated \(F\)-algebra
is detected by its \(F\)-algebra maps to \(F\): every nonzero
element is nonzero under at least one of them.
\end{lemma}

\begin{proof}
Extend \(t\) to a transcendence basis \(t,u_1,\dots,u_j\) of
\(\operatorname{Frac}(B)/F\) by choosing some of the generators of
\(B\). Each generator satisfies a monic polynomial over
\(F(t,u_1,\dots,u_j)\). Choose a single nonzero polynomial \(d(t,u)\)
clearing all the denominators in their coefficients. Then
\(B[d^{-1}]\) is integral over
\(D=F[t,u_1,\dots,u_j,d^{-1}]\).
For all but finitely many \(a\in F\), \(d(a,u)\ne0\):
view the coefficients in the \(u_i\) as polynomials in \(t\) and take
one nonzero coefficient. At such \(a\), specialization at \(t=a\)
exhibits \(t-a\) as a nonunit in \(D\). It is still a nonunit in
\(B[d^{-1}]\): an inverse integral over \(D\) would belong to \(D\),
by multiplying its monic relation by a suitable power of \(t-a\).
It follows that \(t-a\) is a nonunit in \(B\).

A finitely generated \(F\)-algebra that is a field can therefore
have no transcendental element, and finitely many algebraic
generators give a finite extension. Over algebraically closed \(F\),
a maximal ideal in a nonzero finitely generated \(F\)-algebra has
residue field \(F\). To detect a nonzero element \(b\) in a reduced
algebra, apply this to the nonzero finitely generated algebra
obtained by inverting \(b\), and restrict the resulting map.
\end{proof}

\begin{proposition}[No new constants in a simple finite-type algebra]
\label{res:PV05}
Let \(S\) be a simple operator ring finitely generated as an algebra
over \(k\). Then \(S^{\mathrm{const}}=C\); the same holds for
its total ring of fractions.
\end{proposition}

\begin{proof}
By Lemma~\ref{res:PV03}, the simple ring is a finite product of
domains, with one factor in the differential case. Suppose a
constant \(c\in S\) is transcendental over \(k\). Some component
of \(c\) is transcendental, since otherwise the product of
polynomial relations in the finitely many components would give a
relation in \(S\). For every \(a\in C\) the difference \(c-a\) is a
nonzero constant, hence a unit in \(S\). Its component in the
chosen domain is a unit as well. As \(C\) is infinite, this
contradicts Lemma~\ref{res:PV04} with \(F=k\).

Thus \(c\) has a monic relation of least degree in \(k[T]\).
In the differential case, applying \(\partial\) to that relation
gives a relation of lower degree with coefficients the derivatives of
the original coefficients. In the difference case, applying
\(\sigma\) and subtracting the original relation gives one whose
coefficients are the differences of the coefficients and their
images. In both cases the lower-degree relation is zero, hence the
original coefficients lie in \(C\). Since \(S^{\mathrm{const}}\)
is a field and \(C\) is algebraically closed, \(c\) belongs to
\(C\). The statement about the total ring follows from
Lemma~\ref{res:PV03}.
\end{proof}

\section{Construction and uniqueness}

\begin{definition}[Picard--Vessiot ring and its total extension]
\label{res:PV06}
For either system of Lemma~\ref{res:PV02} a Picard--Vessiot ring
over \(k\) is a simple operator extension \(R\) containing a
fundamental matrix \(Y\) and generated as a \(k\)-algebra by its
entries and \(\det(Y)^{-1}\). Its total extension \(Q(R)\) is the
total ring of fractions with the extended operator. It is called a
Picard--Vessiot field when it is a field.
\end{definition}

\begin{theorem}[Existence with prescribed constants]
\label{res:PV07}
Every system of Lemma~\ref{res:PV02} over \(k\) has a
Picard--Vessiot ring. Both the ring and its total extension have
exactly the constants \(C\).
\end{theorem}

\begin{proof}
Let \(U=k[X_{ij},\det(X)^{-1}]\) with \(X\) an indeterminate
square matrix. A derivation of the polynomial ring is defined by
\(\partial X=AX\) and extends to the localization by the
quotient rule. In the difference case set \(\sigma(X)=AX\) and
apply \(\sigma\) to the coefficients. The image of \(\det X\) is
\(\det A\det X\), a unit in \(U\). This is an automorphism
with inverse acting on coefficients by \(\sigma^{-1}\) and on \(X\)
by \(\sigma^{-1}(A^{-1})X\).
By the ascending chain condition, a maximal proper operator
ideal exists, and its quotient is simple. A proper ideal does
not meet \(k\) nontrivially, so the quotient is an extension
of \(k\). The image of \(X\) is fundamental and generates the
required algebra. Proposition~\ref{res:PV05} supplies its constants.
\end{proof}

\begin{lemma}[Descent of ideals with constant variables]
\label{res:PV08}
Let \(S\) be a simple operator ring with constant field \(C\).
If \(D\) is any \(C\)-algebra, equip \(S\otimes_C D\) with the
operator from the first factor only. Every operator ideal is extended
from its contraction to \(D\), and the constants equal \(1\otimes D\).
In particular, if an operator extension \(B\) of \(S\) is generated
as an \(S\)-algebra by a \(C\)-subalgebra \(D\subseteq
B^{\mathrm{const}}\), multiplication gives an isomorphism
\(S\otimes_C D\cong B\).
\end{lemma}

\begin{proof}
For an operator ideal \(I\), divide \(D\) by its contraction and
replace \(I\) by its image. Exactness of tensoring over \(C\) gives
zero contraction for this image. Suppose the image contains a
nonzero \(\sum_{i=1}^m s_i\otimes d_i\) with \(m\) minimal,
and choose the \(d_i\) linearly independent over \(C\).
The possible first coefficients in elements of \(I\) whose
expressions use only \(d_1,\dots,d_m\) form a nonzero operator
ideal of \(S\). It is the whole ring. Thus we can choose such
an element with first coefficient \(1\). Applying \(\partial\),
or subtracting the element from its image under \(\sigma\), now
gives an element of \(I\) with at most \(m-1\) terms. By minimality
it is zero, so all the coefficients are constants. This gives
a nonzero element of the contraction, a contradiction.
Expanding an arbitrary element of the tensor product in a finite
linearly independent set in \(D\) also shows directly that a
constant has constant coefficients. The multiplication map in
the last assertion is surjective with an operator kernel
whose contraction to \(D\) is zero. The first assertion gives
injectivity.
\end{proof}

\begin{theorem}[Uniqueness and recognition inside a simple extension]
\label{res:PV09}
Any two Picard--Vessiot rings for the same system over \(k\) are
isomorphic over \(k\) as operator rings, and the isomorphism
extends to their total extensions. More generally, if a simple
operator extension \(S/k\) has constants \(C\) and contains a
fundamental matrix \(Y\), the subalgebra
\(k[Y_{ij},\det(Y)^{-1}]\) is a Picard--Vessiot ring.
\end{theorem}

\begin{proof}
Take a Picard--Vessiot ring \(R'\) with fundamental matrix \(Y'\).
In \(S\otimes_k R'\) put
\(Z=(Y\otimes1)^{-1}(1\otimes Y')\). Its entries and inverse
determinant are constant by Lemma~\ref{res:PV02}.
The tensor product is nonzero, contains the first factor faithfully,
and is generated over that factor by
\(D=C[Z_{ij},\det(Z)^{-1}]\). Apply Lemma~\ref{res:PV08}
and choose a \(C\)-algebra map \(D\to C\) by
Lemma~\ref{res:PV04}. Evaluation defines an operator map
\(R'\to S\) taking \(Y'\) to \(YB\) for some \(B\in GL_n(C)\).
It is injective by simplicity of \(R'\), and its image is
exactly the algebra generated by \(Y\) and its inverse
determinant, since \(B,B^{-1}\) have entries in \(k\).
This proves recognition. Taking \(S\) itself to be a second
Picard--Vessiot ring gives uniqueness, and localization gives
the assertion about the total extensions.
\end{proof}

\begin{example}[A difference Picard--Vessiot ring that is not a field]
\label{res:PV10}
Let \(k=C(z)\), where \(C\) is algebraically closed of
characteristic zero, and set
\(\sigma(f)(z)=f(qz)\) for a fixed \(q\in C^\times\) of infinite
multiplicative order. The constants are \(C\). For
\(\sigma(y)=-y\), a Picard--Vessiot ring is
\(k[u]/(u^2-1)\) with \(\sigma(u)=-u\). It is already its
own total ring of fractions and has no embedding into a field.
\end{example}

\begin{proof}
A rational function fixed by the substitution \(z\mapsto qz\)
has its zeros and poles invariant under multiplication by \(q\).
Any orbit of a point in \(C^\times\) is infinite, so a nonzero
fixed rational function is \(az^m\) with \(a\in C^\times\) and
\(m\in\mathbb Z\). Invariance gives \(q^m=1\), so \(m=0\).
The idempotents \((1+u)/2,(1-u)/2\) exhibit the asserted ring
as \(k\times k\). The automorphism interchanges the two
idempotents, so no proper nonzero ideal is stable. The
invertible element \(u\) is fundamental and generates the ring,
which is therefore Picard--Vessiot. Its two orthogonal nonzero
idempotents preclude an embedding into a field.
\end{proof}

\paragraph{Sources.}
The differential construction follows the subject of
van der Put--Singer~\cite[Chapter 1, Sections 1.1--1.3]{vps-differential}
and Singer~\cite[Sections 1.2--1.3]{singer}. For the difference
construction see van der Put--Singer~\cite[Chapter 1, Section 1.1]{vps-difference}
and Roques~\cite[Definitions 7--8, Propositions 9 and 12, Theorems 13 and 17]{roques}.
The operator-ring viewpoint also connects to the treatment of Ore
operators in~\cite[Chapter 4]{mg-rm03d}.
